\documentclass[11pt]{article}

\usepackage[a4paper,margin=2.4cm]{geometry}
\usepackage{amsmath,amssymb}
\usepackage{mathtools}
\usepackage{fontspec}
\setmainfont{Liberation Serif}
\usepackage{titlesec}
\usepackage{enumitem}
\usepackage{booktabs}
\usepackage{graphicx}
\usepackage{xcolor}
\usepackage{hyperref}
\hypersetup{colorlinks=true, linkcolor=blue!50!black, urlcolor=blue!50!black}

\title{\textbf{T6C\_fuselage\_cpp} \\ \Large Fuselage Aerodynamic Coefficients --- Physics Formulas by Computation Flow}
\author{}
\date{}

\begin{document}
\maketitle
\vspace{-1.5em}

\noindent This document lists every physics formula in the fuselage
aerodynamic model, in the exact order they are evaluated:
\textbf{configuration \& derived geometry} $\to$
\textbf{flow state} $\to$
\textbf{drag} $\to$
\textbf{pitching moment} $\to$
\textbf{yawing moment} $\to$
\textbf{calibration workflow for when CFD data becomes available}.

\medskip
\noindent \textit{Source}: F.~Nicolosi, P.~Della Vecchia, D.~Ciliberti,
V.~Cusati, ``Fuselage aerodynamic prediction methods,'' \emph{Aerospace
Science and Technology}, 55 (2016) 332--343.

\medskip
\noindent \textit{Why drag and the two moments use different formula
families}: Nicolosi's own method computes $C_D$ from shape factors
$K_n,K_c,K_t$ and computes $C_{M0},C_{M\alpha},C_{N\beta}$ from an
$FR$-term plus nose/tail corrections -- but in the paper \textbf{every
one of those 12 curves is digitized from the authors' own STAR-CCM+ CFD
sweep} (their Figs.~4--6, 8--16), which this project does not have.
The two families differ in how much can be salvaged without it:
\begin{itemize}[nosep]
\item $K_n,K_c,K_t$: \textbf{no} printed numeric value exists anywhere
in the paper's text for these three -- chart-only, zero real anchor.
Rather than invent one from nothing, drag instead uses the closed-form
formulas the paper itself cites as its semi-empirical comparison
baseline (Roskam/Kroo, its Eq.~2,6,7,9) -- real, validated, and
requiring \emph{no} fabricated constant at all (Section~\ref{sec:drag}).
\item $C_{M0},C_{M\alpha},C_{N\beta}$: each \textbf{does} have one real
anchor point -- the paper's reference-fuselage worked example is
printed as actual numbers, not a chart reading (Eq.~13/14/16). No
independently-validated closed-form alternative exists for these (see
project README: Roskam's own fuselage-$C_m$ treatment turns out to be
the same Multhopp lineage, and that lineage's own literature flags it
as unreliable without a case-specific correction factor). So each curve
here is a line through that one real point, with an
\textbf{uncalibrated} trend-guessed slope (Sections~\ref{sec:pitch}--\ref{sec:yaw}).
\end{itemize}

\section{Configuration load (once, at startup)}

Raw geometry input (SI units: metres, radians):
\begin{equation}
L_f,\ d_f,\ L_n,\ L_c,\ L_t,\ d_b,\ \theta,\ \psi,\ S_w,\ \bar{c},\ b
\end{equation}
($\bar c$ = wing mean aerodynamic chord, $b$ = wingspan -- used only for
the optional $S_w$-referenced rescale in Sections~\ref{sec:cdassembly},
\ref{sec:pitchassembly}, \ref{sec:yawassembly}).

Derived fineness ratios:
\begin{equation}
FR = \frac{L_f}{d_f}, \qquad FR_n = \frac{L_n}{d_f}, \qquad FR_t = \frac{L_t}{d_f}
\end{equation}

Frontal area and wetted area (exact nose-cone / cabin-cylinder /
tailcone-frustum decomposition -- closed-form solid-of-revolution
geometry, not an empirical fit):
\begin{equation}
S_{\text{front}} = \frac{\pi}{4}d_f^{\,2}
\end{equation}
\begin{align}
S_{\text{wet,nose}} &= \pi\,\frac{d_f}{2}\,\ell_n, & \ell_n &= \sqrt{L_n^2+\left(\tfrac{d_f}{2}\right)^2} \\
S_{\text{wet,cabin}} &= \pi\,d_f\,L_c \\
S_{\text{wet,tail}} &= \pi\left(\tfrac{d_f}{2}+\tfrac{d_b}{2}\right)\ell_t, & \ell_t &= \sqrt{L_t^2+\left(\tfrac{d_f-d_b}{2}\right)^2} \\
S_{\text{wet}} &= S_{\text{wet,nose}}+S_{\text{wet,cabin}}+S_{\text{wet,tail}}
\end{align}

$FR,FR_n,FR_t,S_{\text{front}},S_{\text{wet},\bullet}$ depend only on
fixed geometry, so they are computed once in the constructor and
cached -- \texttt{Update()} never recomputes them.

\section{Flow state (every call)}

\begin{equation}
q=\tfrac12\rho V^2, \qquad Re=\frac{\rho V L_f}{\mu}
\end{equation}
\begin{equation}
C_{Dfp} = \frac{0.455}{\left(\log_{10}Re\right)^{2.58}\left(1+0.144M^2\right)^{0.65}}
\qquad\text{(Nicolosi Eq.~2)}
\end{equation}

\section{Fuselage drag --- Roskam/Kroo closed-form}
\label{sec:drag}

Skin friction (Eq.~6; the extracted-PDF text read the fineness-ratio
coefficient as $0.0225$, but that produces $C_{Dsf}$ alone exceeding the
paper's own reference-case DATCOM total -- $0.0025$ (matching the
well-known Raymer/Hoerner fuselage form factor $(l/d)/400$) reproduces
the paper's Table~3 reference value to within 2.4\%, see
\texttt{main.cpp}'s regression test):
\begin{equation}
C_{Dsf} = C_{Dfp}\left[1+\frac{60}{(L_f/d_f)^3}+0.0025\frac{L_f}{d_f}\right]\frac{S_{\text{wet}}}{S_{\text{front}}}
\end{equation}
Base drag (Eq.~7, zero if $d_b=0$):
\begin{equation}
C_{Dbase} = \frac{0.029(d_b/d_f)^3}{\sqrt{C_{Dsf}\,S_{\text{front}}/S_{\text{wet}}}}\,\frac{S_{\text{wet}}}{S_{\text{front}}}
\end{equation}
Upsweep (Eq.~9, Kroo; the extracted-PDF text read the area ratio as
$S_w/S_{\text{front}}$, but that produces $C_{Dupsweep}$ alone over
double the paper's own reference-case DATCOM total -- inverting it to
$S_{\text{front}}/S_w$ reproduces Table~3 to within 2.4\% together with
the corrected $C_{Dsf}$ above; linearized-centreline approximation for
$(h/l)_{0.75lt}$, this project's own addition -- Kroo's formula needs
the tailcone centreline's vertical rise at 75\% of tailcone length,
non-dimensionalized by tailcone length; approximating the rise as
linear from cabin to tail tip makes the ratio independent of $L_t$):
\begin{equation}
\left(\frac hl\right)_{0.75lt}\approx0.75\tan\theta,
\qquad
C_{Dupsweep} = 0.075\,\frac{S_{\text{front}}}{S_w}\,(0.75\tan\theta)
\end{equation}
Windshield contribution (Eq.~8) omitted -- needs a graphical Roskam
chart, no closed form; a calibration hook is added on top instead:
\begin{equation}
C_{D,\text{fus}}^{(S_{\text{front}})} = C_{Dsf}+C_{Dbase}+C_{Dupsweep} + \texttt{windshield\_drag\_increment} \quad(\text{default }0)
\end{equation}

\section{Drag assembly}
\label{sec:cdassembly}

Rescaled to the wing-area reference used by \texttt{CDo}/\texttt{CDgear}
in \texttt{t6texan2.xml}:
\begin{equation}
C_{D,\text{fus}}^{(S_w)} = C_{D,\text{fus}}^{(S_{\text{front}})}\,\frac{S_{\text{front}}}{S_w},
\qquad
D_{\text{fus}} = q\,S_{\text{front}}\,C_{D,\text{fus}}^{(S_{\text{front}})}
\end{equation}

\section{Zero-angle-of-attack pitching moment $C_{M0}$}
\label{sec:pitch}

Structure (Nicolosi Eq.~11):
\begin{equation}
C_{M0,\text{fus}} = C_{M0,FR}(FR) + \Delta C_{M0,\text{nose}}(FR_n,\psi) + \Delta C_{M0,\text{tail}}(FR_t,\theta)
\end{equation}
Each term is a line through the \textbf{one real anchor point the paper
prints as numbers} (Eq.~13, reference fuselage $FR=8.7$, $FR_n=1.6$,
$\psi=40^\circ$, $FR_t=2.8$, $\theta=14^\circ$), not a chart reading:
\begin{align}
C_{M0,FR}(FR) &= \underbrace{-0.3159}_{\text{Eq.13 anchor}} + s_{FR}\,(FR-8.7) \\
\Delta C_{M0,\text{nose}}(FR_n,\psi) &= \underbrace{0.0013}_{\text{Eq.13 anchor}} + s_{n,\psi}(\psi-40^\circ) + s_{n,FR_n}(FR_n-1.6) \\
\Delta C_{M0,\text{tail}}(FR_t,\theta) &= \underbrace{0.0013}_{\text{Eq.13 anchor}} + s_{t,\theta}(\theta-14^\circ) + s_{t,FR_t}(FR_t-2.8)
\end{align}
Slopes $s_{FR},s_{n,\psi},s_{n,FR_n},s_{t,\theta},s_{t,FR_t}$ are
\texttt{TuningParams::pitch} fields, defaulted small (paper: nose
correction $\sim15\%$ of $C_{M0,FR}$, tail correction $<10\%$,
Sec.~5.1) and \textbf{uncalibrated}.

\section{Pitching moment derivative $C_{M\alpha}$}

Same structure (Nicolosi Eq.~12), anchored at Eq.~14:
\begin{align}
C_{M\alpha,FR}(FR) &= \underbrace{0.1815}_{\text{Eq.14 anchor}} + s_{\alpha,FR}\,(FR-8.7) \\
\Delta C_{M\alpha,\text{nose}}(FR_n,\psi) &= \underbrace{0.0005}_{\text{Eq.14 anchor}} + \ldots \\
\Delta C_{M\alpha,\text{tail}}(FR_t,\theta) &= \underbrace{0.0006}_{\text{Eq.14 anchor}} + \ldots
\end{align}
(paper: these two corrections change ``about 1\% or less'' with moment
reference point, Sec.~5.2 -- default slopes accordingly small).
\begin{equation}
C_{M\alpha,\text{fus}} = C_{M\alpha,FR}+\Delta C_{M\alpha,\text{nose}}+\Delta C_{M\alpha,\text{tail}}
\end{equation}

\section{Pitching moment assembly}
\label{sec:pitchassembly}

\begin{equation}
C_{M,\text{fus}} = C_{M0,\text{fus}} + C_{M\alpha,\text{fus}}\,\alpha
\qquad\text{(Nicolosi Eq.~10)}
\end{equation}
Reference is $S_{\text{front}}$ and $d_f$ (paper's own convention):
\begin{equation}
M_{\text{fus}} = q\,S_{\text{front}}\,d_f\,C_{M,\text{fus}}
\end{equation}
Optional rescale to the wing-area/$\bar c$ reference (paper's Tables~6--7 form):
\begin{equation}
C_{M,\text{fus}}^{(S_w,\bar c)} = C_{M,\text{fus}}\,\frac{S_{\text{front}}\,d_f}{S_w\,\bar c}
\end{equation}

\section{Yawing moment derivative $C_{N\beta}$}
\label{sec:yaw}

Structure (Nicolosi Eq.~15), anchored at Eq.~16 -- note only $FR_n$/$FR_t$
appear (no $\psi,\theta$ dependence stated for yaw, Sec.~6):
\begin{align}
C_{N\beta,FR}(FR) &= \underbrace{0.1817}_{\text{Eq.16 anchor}} + s_{\beta,FR}\,(FR-8.7) \\
\Delta C_{N\beta,\text{nose}}(FR_n) &= \underbrace{-0.0003}_{\text{Eq.16 anchor}} + s_{\beta n}\,(FR_n-1.6) \\
\Delta C_{N\beta,\text{tail}}(FR_t) &= \underbrace{-0.0002}_{\text{Eq.16 anchor}} + s_{\beta t}\,(FR_t-2.8)
\end{align}
\begin{equation}
C_{N\beta,\text{fus}} = C_{N\beta,FR}+\Delta C_{N\beta,\text{nose}}+\Delta C_{N\beta,\text{tail}}
\end{equation}

\section{Yawing moment assembly}
\label{sec:yawassembly}

\begin{equation}
C_{N,\text{fus}} = C_{N\beta,\text{fus}}\,\beta,
\qquad
N_{\text{fus}} = q\,S_{\text{front}}\,d_f\,C_{N,\text{fus}}
\end{equation}
Optional rescale to the wing-area/span reference, and the scaling the
paper itself specifies if this feeds a whole-aircraft $C_{N\beta}$
(Sec.~6, before adding the vertical tail's contribution):
\begin{equation}
C_{N,\text{fus}}^{(S_w,b)} = C_{N,\text{fus}}\,\frac{S_{\text{front}}\,d_f}{S_w\,b}
\end{equation}

\section{Direction, deliberately left to the caller}
$D_{\text{fus}}$ is a scalar magnitude opposing the relative wind;
$M_{\text{fus}},N_{\text{fus}}$ are scalar moments about the paper's
reference point ($x/L_f{=}0.50$, $z/d_f{=}0.50$). None of the three are
resolved into body axes here -- the caller rotates/relocates them the
same way it treats any other aerodynamic contribution.

\section{Calibration workflow --- filling in real CFD data later}
\label{sec:calibration}

$C_D$ (Section~\ref{sec:drag}) needs no calibration -- Roskam/Kroo is
closed-form already. Only $C_{M0},C_{M\alpha},C_{N\beta}$
(Sections~\ref{sec:pitch}--\ref{sec:yaw}) carry uncalibrated slopes.
When a real CFD run (STAR-CCM+, OpenFOAM, whatever) on an actual T-6C
fuselage geometry becomes available, replace them with fitted ones --
\textbf{never edit the formulas in \texttt{t6c\_fuselage.cpp}, only the
numbers in \texttt{t6c\_fuselage\_params.yaml}}:

\begin{enumerate}[nosep]
\item Run at least two geometries through CFD (e.g.\ two different $FR$
at fixed nose/tail, then two different $FR_n$ at fixed $FR/FR_t$, etc.)
and extract $C_{M0},C_{M\alpha},C_{N\beta}$ for the whole fuselage at
each, the same way Nicolosi generated Figs.~8--16 from many CFD cases.
\item Fit a real slope $s_{FR}$ (etc.) through those points instead of
the trend-guess default -- with only 2 points this is just
$(y_2-y_1)/(x_2-x_1)$; with more, a least-squares line.
\item With $\geq3$--4 real data points per curve, consider replacing
the linear trend-guess with the generic
\texttt{ChartTable1D}/\texttt{ChartTable2D} piecewise-linear
interpolation described as future work in the project README, instead
of the fixed linear form used here.
\item Update \texttt{t6c\_fuselage\_params.yaml} only; re-run
\texttt{main.cpp}'s regression test to confirm the new numbers still
reproduce whatever CFD points were used to fit them.
\end{enumerate}

\bigskip
\noindent\textbf{Call graph, one call per evaluation} (no persistent
state):
\begin{quote}
\texttt{FuselageModel::FuselageModel(geometry, tuning)} --- Section~1,
once, at construction\\
\texttt{FuselageModel::Update(state)} --- Sections~2--9, every call,
pure function of \texttt{state} (which carries $\alpha,\beta$ in
addition to $V,\rho,\mu,M$) and the cached Section~1 geometry
\end{quote}

\end{document}
